\section{An equivalent axiom emitter with 39 fewer source instructions}
\label{sec:distributed-gate}
\begingroup
\newcommand{\Imp}{\mathsf I}
\newcommand{\Mem}{\mathsf M}
\newcommand{\All}{\mathsf U}
\newcommand{\Eq}{\mathsf E}
\newcommand{\Neg}{\mathsf N}

A new source-level simplification distributes the innermost universal
quantifier over the two conjuncts of PA196's separation condition. The
result has \textbf{199 AX macro tokens instead of 201}, and expands to
\textbf{3,773 source instructions instead of 3,812}. These are source
counts. They establish no reduction in physical Turing-machine states.

Use the exact PA196 constructors
\[
 \Imp(x,y)=4p(x,y),\quad
 \Mem(x,y)=4p(x,y)+1,\quad
 \All(x,B)=4p(B,x)+2,\quad
 \Eq(x,y)=4p(x,y)+3,
\]
where $p$ is the upstream Cantor pairing, and put
$\Neg(B)=\Imp(B,1)$ and
$B\wedge C=\Neg(\Imp(B,\Neg(C)))$.
The property parameter $c$ is simultaneously a formula code and the
fresh-output variable label, exactly as in the upstream emitter.

For the original gate parameters $(a,b,c)$ define
\begin{align*}
 \Delta&=\Eq(a,b)\to\All(c,\Eq(a,b)), &
 V&=\Mem(a,b)\to\All(c,\Mem(a,b)),\\
 P&=c\to\Mem(c,b), &
 A&=\Mem(a,b)\to P, &
 B&=\Mem(c,b)\to c.
\end{align*}
Here an arrow denotes $\Imp$, and $c$ in formula position denotes the
property with numeric code $c$. The original definitions are
\[
 S=\All(b,A\wedge B),\qquad
 Q=\All(c,S\to\Neg(P)),\qquad
 F=\Neg\bigl(\Delta\to(V\to Q)\bigr).
\]
Replace them by
\begin{align*}
 S_1&=\All(b,A), & S_2&=\All(b,B),\\
 Q_*&=\All\bigl(c,S_1\to(S_2\to\Neg(P))\bigr), &
 F_*&=\Neg\bigl(\Delta\to(V\to Q_*)\bigr).
\end{align*}
The guard remains $c\ne a,b$; its disabled result remains truth code zero.
The other six heads, the restricted propositional head, and the saved
incoming/generalization entries are unchanged. In particular, the AX
command still pushes nine items, and the head-selection ranks and
fixed push/pop words require no change merely because of this logical
reformulation. The numerical value of the enabled gate is different;
the claim is derivable logical equivalence, not equality of codes.

\begin{proposition}[Gate equivalence in the common numeric calculus]
\label{prop:distributed-gate}
The retained heads h1, h6, the restricted implicational basis, modus
ponens, and generalization derive both $F\to F_*$ and $F_*\to F$ for
all numeric parameters. The proof does not use either version of the
guarded gate, equality symmetry, a substitution rule, or unrestricted
universal instantiation.
\end{proposition}
\begin{proof}
The unchanged restricted implicational basis yields the classical
implicational rules. With unchanged h1 it also yields explosion for
code~1, hence the usual rules for the displayed negative encoding of
conjunction. These are the upstream bootstrap arguments, independent of
the gate.

For arbitrary formula codes $X,Y$ and label $b$, generalize the
propositional projections $(X\wedge Y)\to X$ and
$(X\wedge Y)\to Y$. Applying h6 gives
\[
 \All(b,X\wedge Y)\to\All(b,X),\qquad
 \All(b,X\wedge Y)\to\All(b,Y).
\]
Conversely, generalize the propositional theorem
$X\to(Y\to X\wedge Y)$ and apply h6 twice. This gives
\[
 \All(b,X)\to
 \bigl(\All(b,Y)\to\All(b,X\wedge Y)\bigr).
\]
Propositional reasoning therefore proves
\[
 \All(b,X\wedge Y)\ \leftrightarrow\
       \bigl(\All(b,X)\wedge\All(b,Y)\bigr).
\]
This is normal-quantifier distribution over a binary conjunction. It
requires no freshness side condition.

Apply it with $X=A$ and $Y=B$. Currying the two conjuncts now gives
\[
 S\to\Neg(P)
 \quad\leftrightarrow\quad
 S_1\to(S_2\to\Neg(P)).
\]
Generalize each direction at $c$ and use h6 to obtain
$Q\leftrightarrow Q_*$. Implicational reasoning transports this
equivalence through the unchanged contexts
$\Neg(\Delta\to(V\to\cdot))$, proving the result.
\end{proof}

The preceding proof is about the actual numerical constructors, not
about a change of variable names in a decoded first-order language.
The inner label $b$ binds exactly the same occurrences on both sides
of the conjunction-distribution law. If $a=b$, the resulting additional
bound occurrences are present on both sides. The outer label $c$ still
binds the whole curried condition, and the occurrences of $\Delta$ and
$V$ remain outside it. The simultaneous use of $c$ as a property code
and a bound label is preserved. Thus no freshness assumption beyond
the original gate guard enters, and no general substitution theorem
is being used implicitly.

\subsection{Uniform proof translation and its PA verification}

Let $C$ be the exact numerical proof calculus with the original gate,
and let $C_*$ replace only its gate by $F_*$. A proof is a finite list
of labelled lines. An axiom line records its head and numerical
parameters; an MP line records the two earlier line indices and its
conclusion; a generalization line records an earlier line and its
label. The checker verifies numerical equalities of the displayed
constructors, the restriction $u\leq a$ on the propositional head,
and the original two inequalities on an enabled gate. One may code
lists, for example, by $\operatorname{nil}=0$ and
$\operatorname{cons}(h,t)=1+p(h,t)$. Use any fixed tagged coding of
these line records. Write $\operatorname{Prf}_C(p,f)$ and
$\operatorname{Prf}_{C_*}(p,f)$ for the resulting primitive recursive
proof predicates, with last conclusion $f$.

\begin{theorem}[PA verification of the numerical proof translation]
\label{thm:distributed-pa-translation}
There are primitive recursive proof translators $\tau$ and $\sigma$
for which PA proves
\begin{align*}
 \operatorname{Prf}_C(p,f)&\longrightarrow
   \operatorname{Prf}_{C_*}(\tau(p),f),\\
 \operatorname{Prf}_{C_*}(p,f)&\longrightarrow
   \operatorname{Prf}_{C}(\sigma(p),f).
\end{align*}
Consequently PA proves equality of the two refutability predicates,
\[
 \exists p\,\operatorname{Prf}_C(p,1)
 \quad\longleftrightarrow\quad
 \exists p\,\operatorname{Prf}_{C_*}(p,1),
\]
and hence equality of their consistency assertions.
\end{theorem}
\begin{proof}
The restricted implicational bootstrap is a fixed finite family of
proof templates. Its constructor bounds follow from
$x,y\leq p(x,y)\leq4p(x,y)$. In particular it supplies the usual
$K,S$, Peirce, composition, and exchange templates at arbitrary
numerical inputs. The pinned upstream files
\file{M6/Logic/RestrictedImplication.lean} and
\file{PA196/Logic/ImplicationBasis.lean} give these explicit finite
constructions. They use no gate. The finite-context deduction
construction in \file{PA196/Logic/Deduction.lean} uses only MP and
closed theorems, so it never generalizes an undischarged assumption.

For completeness, explosion also has a uniform gate-free template.
From a temporary assumption $1$, the $K$ template gives
\[
 \All(0,\Neg(\All(0,X)))\to1
       =\exists 0\,\All(0,X).
\]
Head h1 then gives $X$; the finite deduction construction closes this
as $1\to X$. Peirce and explosion supply double-negation elimination,
and the three conjunction templates used in
Proposition~\ref{prop:distributed-gate} follow by ordinary finite
deduction. The upstream \file{PA196/Logic/Connectives.lean} contains
those constructors explicitly. The only generalizations in the gate
equivalence proof are of already closed conjunction or implication
theorems. Thus each direction of that proof expands to a fixed finite
template whose parameters are $a,b,c$.

Scan an input proof from left to right, maintaining the index of the
translated conclusion of each old line. A changed gate-axiom line
with conclusion $F(a,b,c)$ is replaced by the $C_*$ gate axiom
$F_*(a,b,c)$, the common-core template proving $F_*\to F$, and MP.
Unchanged axioms are copied. An MP or generalization line is copied
with its prior-line references updated. Every translated line group
therefore ends in exactly the original numerical conclusion code;
there is no formula recoding. The reverse translator uses $F\to F_*$.
Extend both algorithms arbitrarily to malformed input records, for
example by outputting the empty proof when a required record is
missing. List decoding, constructor calculation, finite-template
instantiation, appending, and this bounded scan are primitive
recursive.

PA verifies each fixed template by its finitely many checking
equations and the elementary pairing bounds above. Induction on
the input proof length proves the invariant that the remembered
output index has the same numerical conclusion and a valid proof
prefix. The only changed axiom has exactly the same numerical side
conditions in the two calculi. This proves the two displayed
implications inside PA. Existential introduction in each direction
proves the final equivalence. This is a syntactic proof-checking
argument, and does not use soundness, reflection, or a truth predicate.
\end{proof}

The code-zero identity $\Imp(0,0)=0$ does not obstruct this argument:
the translation preserves all original conclusion codes, including
zero. If the source-level checker exposes a disabled gate as code
zero, that result is identical on both sides; it also has the common
identity proof $0\to0$. The theorem is an ordinary mathematical
arithmetization proof. The new translators have not been exported as
Lean declarations in this artifact.

\paragraph{The selected head order.}
The grammar candidate used by the physical compiler additionally
emits the six heads in order $[1,6,2,3,4,5]$, and writes its sixth head
in the exchanged form
\[
 \All(D,A\to H)\to\bigl(\All(D,A)\to\All(D,H)\bigr).
\]
The common restricted implicational basis proves exchange independently
of h6. Thus the original and exchanged h6 derive one another by one
exchange-template application and MP. Add that fixed replacement
case to $\tau$ and $\sigma$ to obtain
Theorem~\ref{thm:distributed-pa-translation} for this selected
calculus as well. Permuting the emitted heads does not change the
axiom family, but it changes the stack rank used to select each
head. The compiler must use the corresponding new ranks. The
independent word audit checks this precise permutation and the
precise exchanged head; it does not treat an arbitrary permutation
as an unchanged numerical command.

\paragraph{Additional equivalent words.}
Inside $Q_*$, the three antecedents $S_1,S_2,P$ precede the final
falsity code. All six permutations are equivalent by repeated
implicational exchange. Within $S_1$, the two antecedents of
$\Mem(a,b)\to(c\to\Mem(c,b))$ may also be exchanged.
These choices give twelve explicitly certified words with the same
macro multiset. They change formula-code layouts and can therefore
change grammar sharing even though they prove the same formulas.

\subsection{Exact expansion count}

The accompanying \file{distributed_gate.py} reconstructs the original
201-token word from the mathematical heads and compares it with the
literal blocks in upstream \file{PA196/Source/AXProgram.lean}. It checks
the pinned source identity and contiguous block boundaries before
emitting any variants. The new words each have 199 tokens. Their
multiset difference from the original is
\[
  +1\;\texttt{copyD}\ -1\;\texttt{cons}\ -2\;\texttt{push}.
\]
All increment counts remain the same. The authoritative block lengths
are 26 primitive instructions for \texttt{copyD}, 27 for
\texttt{cons}, and 19 for \texttt{push}. Consequently
\[
  26-27-2\cdot19=-39,
  \qquad 3812-39=3773.
\]
The enabled gate's contribution drops from 77 to 75 macro tokens.
The simplification replaces the negative conjunction's two constant-one
constructions and three pairing operations by a second universal
constructor and a curried antecedent. The source arithmetic above
accounts for their different expansion lengths.

The source is pinned to
\nolinkurl{jamesyc/BusyBeaver@407cdc356f6b586635e71ac92be394255f4bfea4}.
The checked SHA-256 of \file{PA196/Program.lean} is
\begin{quote}\small
\nolinkurl{8c7c6c6884ecbb8f56e893404bb81162db15ef3dd6fd3212881a93994d63d872}.
\end{quote}
The AX-block source has SHA-256
\begin{quote}\small
\nolinkurl{8eda8dc0ed0b445c05726990c5a07634464c10ba547fff6bbdba516578ad572c}.
\end{quote}
The receipt records all twelve words, exact counts, and the explicit
absence of an established physical state count.

\paragraph{What this does not justify.}
The equality- and membership-vacuity factors have deliberately remained
outside the outer universal quantifier. Moving them through that
quantifier would require an additional derivation; their reuse of the
same bound label cannot be dismissed by an informal freshness claim.
The present transformation only distributes a universal over a
conjunction and curries antecedents, operations proved above from the
unchanged common heads. A source tree, physical allocation, blank
startup, universal runtime simulation, and the complete PA consistency
calibration remain separate requirements for a smaller TM.

\endgroup
